% corpus_supplement: Erdős problem #617
% source: https://www.erdosproblems.com/latex/617
% fetched_at: 2026-07-14T18:48:19Z
% payload_sha256: fbdadfc6ce07ed82cb552e9af24f88bd336af2653f665e8a59f34af86331e485
% statement/remarks/references extracted by erdos_ingest.extract_source_page;
% rendered by erdos_ingest.render_tex — do not edit
\documentclass{article}
\usepackage{amsmath, amssymb, amsthm}
\usepackage{hyperref}
\usepackage{geometry}
\geometry{margin=1in}

\title{Erd\H{o}s Problem \#617}
\date{Last updated: 2025-08-31}
\author{Source: erdosproblems.com}

\begin{document}
\maketitle

\noindent\textbf{Status:} OPEN \quad \textbf{No prize}

\noindent\textbf{Tags:} graph theory

\medskip
\noindent\textbf{Problem Statement:}

Let $r\geq 3$. If the edges of $K_{r^2+1}$ are $r$-coloured then there exist $r+1$ vertices with at least one colour missing on the edges of the induced $K_{r+1}$.

\bigskip
\noindent\textbf{Remarks:}

In other words, there is no balanced colouring. A conjecture of Erd\H{o}s and Gy\'{a}rf\'{a}s \cite{ErGy99}, who proved it for $r=3$ and $r=4$ (and observed it is false for $r=2$), and showed this property fails for infinitely many $r$ if we replace $r^2+1$ by $r^2$.

\bigskip
\noindent\textbf{References:}

[ErGy99] Erd\H{o}s, Paul and Gy\'{a}rf\'{a}s, Andr\'{a}s, Split and balanced colorings of complete graphs. Discrete Math. (1999), 79-86.

\bigskip
\noindent\small{Source: \url{https://www.erdosproblems.com/617}}
\end{document}
